\documentclass[12pt]{article}
\usepackage[utf8]{inputenc}
\usepackage[T1]{fontenc}
\usepackage[margin=1in]{geometry}
\usepackage{setspace}
\usepackage{xcolor}
\usepackage[colorlinks=true, urlcolor=blue!60!black, linkcolor=blue!60!black]{hyperref}
\setlength{\parindent}{2em}
\setlength{\parskip}{1.25em}
\renewcommand{\baselinestretch}{1.0}
\title{\textit{Balzac v. Porto Rico}, 258 U.S. 298 (1922)}
\author{Case Brief}
\date{}
\begin{document}
\maketitle

\section{Case Metadata}
\begin{itemize}
  \item \textbf{Case Name:} \textit{Balzac v. Porto Rico}, 258 U.S. 298 (1922)
  \item \textbf{Chapter:} Chapter 3: Federal Legislative Power
  \item \textbf{Court and Dates:} Supreme Court of the United States; argued March 20, 1922; decided April 10, 1922
  \item \textbf{Lower Court:} Supreme Court of Porto Rico, \textit{People v. Balzac}, 28 P.R.R. 139 \& 141 (affirmed)
  \item \textbf{Vote:} 9--0 (unanimous judgment)
  \item \textbf{Majority Author:} Chief Justice William Howard Taft (joined by Justices McKenna, Day, Van Devanter, Pitney, McReynolds, Brandeis, and Clarke)
  \item \textbf{Concurrence in the Result:} Justice Oliver Wendell Holmes Jr. (without separate opinion)
  \item \textbf{Primary Texts:} \href{https://supreme.justia.com/cases/federal/us/258/298/}{Justia U.S. Supreme Court Center}; \href{https://en.wikisource.org/wiki/Balzac_v._Porto_Rico/Opinion_of_the_Court}{Wikisource Full Opinion}
  \item \textbf{Constitutional and Statutory Provisions:} U.S. Const. art.~IV, \S~3, cl.~2 (\href{https://www.law.cornell.edu/constitution-conan/article-4/section-3/clause-2/property-clause}{Territories Clause}); art.~III, \S~2; First and Sixth Amendments; Jones Act (Organic Act of Porto Rico of March 2, 1917), 39 Stat. 951, \S\S~2, 5
\end{itemize}

\section{Facts and Procedural History}
Jesus M. Balzac was the editor of \textit{El Baluarte}, a daily newspaper published in Arecibo, Porto Rico (\href{https://supreme.justia.com/cases/federal/us/258/298/}{258 U.S. at 300}). On April 16 and April 23, 1918---just over a year after Congress enacted the Jones Act of March 2, 1917, whose Section~5 conferred collective U.S. citizenship upon citizens of Porto Rico---Balzac published two articles attacking the conduct of the Governor of the island. The district attorney filed two informations charging Balzac with criminal libel in the District Court for Arecibo.

Under Porto Rico's Penal Code (\S\S~13--14, 244), criminal libel was classified as a misdemeanor punishable by up to two years in jail or a \$5,000 fine, and Section~178 of the Code of Criminal Procedure granted jury trials only in felony cases (a subsequent 1919 statute extending jury trials to misdemeanors took effect after Balzac's trials). In both prosecutions, Balzac demanded a jury trial under the Sixth Amendment as an American citizen and argued that his articles were protected fair comment under the First Amendment. The trial court overruled both defenses, tried Balzac without a jury, convicted him in both cases, and sentenced him to five months' imprisonment in the district jail in the first case and four months in the second, plus costs (\href{https://en.wikisource.org/wiki/Balzac_v._Porto_Rico/Opinion_of_the_Court}{258 U.S. at 300}). The Supreme Court of Porto Rico affirmed both convictions, and Balzac brought writs of error to the U.S. Supreme Court.

\section{Issue(s) Presented}
\begin{enumerate}
  \item Did Congress's enactment of the Jones Act of 1917---particularly its grant of U.S. citizenship to Puerto Ricans---incorporate Porto Rico into the United States so that the Sixth Amendment right to a criminal jury trial applied there of its own force?
  \item Were Balzac's newspaper articles protected as fair comment under the First Amendment?
\end{enumerate}

\section{Holding}
No on both counts. Writing for a unanimous Court in the judgment (with Justice Holmes concurring in the result), Chief Justice Taft held that the Jones Act of 1917 did \textit{not} incorporate Porto Rico into the Union. Cementing the \textit{Insular Cases} doctrine under the Territories Clause (art.~IV, \S~3, cl.~2), the Court ruled that only ``fundamental'' constitutional rights apply automatically in unincorporated territories, that the Sixth Amendment right to trial by jury is not such a fundamental right, and that it is ``locality,'' not U.S. citizenship, that governs whether federal jury-trial guarantees apply (\href{https://supreme.justia.com/cases/federal/us/258/298/}{258 U.S. at 304--13}). The Court also summarily rejected Balzac's First Amendment defense, calling his articles ``excessive and outrageous'' (\href{https://supreme.justia.com/cases/federal/us/258/298/}{id. at 314}).

\section{Reasoning}
After confirming writ-of-error jurisdiction under Judicial Code \S\S~237 and 246 (finding that Balzac's Sixth Amendment demand sufficiently challenged the validity of Porto Rico's misdemeanor bench-trial statutes), Chief Justice Taft consolidated the fragmented \textit{Insular Cases} into a single, bright-line constitutional framework:

\begin{enumerate}
  \item \textbf{Adoption of Justice White's Incorporation Doctrine:} Acknowledging earlier ``diversity of opinion'' in the \textit{Insular Cases}, Taft declared that Justice Edward Douglass White's concurring framework in \textit{Downes v. Bidwell}, 182 U.S. 244 (1901), as confirmed in \textit{Dorr v. United States}, 195 U.S. 138 (1904), had become ``the settled law of the court'' (\href{https://en.wikisource.org/wiki/Balzac_v._Porto_Rico/Opinion_of_the_Court}{258 U.S. at 305}). Under that rule, constitutional jury-trial guarantees (art.~III, \S~2; Sixth and Seventh Amendments) apply to \textit{incorporated} territories destined for statehood, but not to \textit{unincorporated} territories ``merely belonging to'' the United States.
  \item \textbf{Clear-Statement Presumption Against Incorporation:} Following the Spanish-American War, congressional intent to incorporate new territory ``is not to be assumed without express declaration, or an implication so strong as to exclude any other view'' (\href{https://supreme.justia.com/cases/federal/us/258/298/}{258 U.S. at 306}). Neither the Jones Act's title nor its text declared incorporation. Crucially, Section~2 enacted a statutory ``Bill of Rights'' mirroring most federal guarantees while deliberately omitting grand and petit juries---proof that Congress did not believe the federal Bill of Rights applied \textit{ex proprio vigore}.
  \item \textbf{Citizenship Without Incorporation---``Locality'' Controls:} Section~5's grant of U.S. citizenship did not imply incorporation. Taft explained that Congress conferred citizenship to satisfy Puerto Ricans' ``yearning'' for international diplomatic protection and to allow them to move to the continental United States and enjoy full political rights there without naturalization (\href{https://en.wikisource.org/wiki/Balzac_v._Porto_Rico/Opinion_of_the_Court}{258 U.S. at 308}). Within Porto Rico, however, ``[i]t is locality that is determinative of the application of the Constitution, in such matters as judicial procedure, and not the status of the people who live in it'' (\href{https://supreme.justia.com/cases/federal/us/258/298/}{id. at 309}). Even a mainland U.S. citizen residing in Porto Rico enjoyed no federal constitutional right to a jury trial there.
  \item \textbf{Distinguishing Alaska and Invoking Cultural/Racialized Assumptions:} Unlike Alaska---which \textit{Rasmussen v. United States}, 197 U.S. 516 (1905), treated as incorporated because it was a contiguous, sparsely settled continental territory open to mainland settlement---Porto Rico was a densely populated, ``distant ocean'' community of ``different origin and language'' trained in Spanish civil law (\href{https://en.wikisource.org/wiki/Balzac_v._Porto_Rico/Opinion_of_the_Court}{258 U.S. at 309--11}). Because jury service was an ``institution of Anglo-Saxon origin,'' Congress properly let local residents decide when to adopt it rather than forcing it on them.
  \item \textbf{Non-Article~III Territorial Courts and Fundamental Rights:} Finally, establishing a U.S. District Court in Porto Rico did not imply incorporation; it was a legislative territorial court created under Article~IV, Section~3, not a ``true United States court'' under Article~III. While ``certain fundamental personal rights'' such as due process always constrained federal power in unincorporated territories, trial by jury was not among them (\href{https://en.wikisource.org/wiki/Balzac_v._Porto_Rico/Opinion_of_the_Court}{258 U.S. at 312--13}).
\end{enumerate}

\section{Concurring and Dissenting Opinions}
No Justice dissented. Justice Oliver Wendell Holmes Jr. concurred solely in the result without issuing an opinion (\href{https://supreme.justia.com/cases/federal/us/258/298/}{258 U.S. at 314}). Scholars have long noted that Holmes's bare concurrence signaled discomfort with Taft's expansive territorial-incorporation rationale (or with the First Amendment dismissal) even though Holmes agreed the convictions should stand.

\section{Key Quotations}
\begin{quote}
``Before that, the purpose of Congress might well be a matter of mere inference from various legislative acts; but in these latter days, incorporation is not to be assumed without express declaration, or an implication so strong as to exclude any other view.'' --- Chief Justice Taft, \href{https://supreme.justia.com/cases/federal/us/258/298/}{258 U.S. at 306}
\end{quote}

\begin{quote}
``It is locality that is determinative of the application of the Constitution, in such matters as judicial procedure, and not the status of the people who live in it.'' --- Chief Justice Taft, \href{https://supreme.justia.com/cases/federal/us/258/298/}{258 U.S. at 309}
\end{quote}

\begin{quote}
``Congress has thought that a people like the Filipinos, or the Porto Ricans, trained to a complete judicial system which knows no juries, living in compact and ancient communities, with definitely formed customs and political conceptions, should be permitted themselves to determine how far they wish to adopt this institution of Anglo-Saxon origin, and when.'' --- Chief Justice Taft, \href{https://en.wikisource.org/wiki/Balzac_v._Porto_Rico/Opinion_of_the_Court}{258 U.S. at 310}
\end{quote}

\begin{quote}
``It is past time to acknowledge the gravity of this error and admit what we know to be true: The Insular Cases have no foundation in the Constitution and rest instead on racial stereotypes. They deserve no place in our law.'' --- Justice Gorsuch (concurring), \href{https://supreme.justia.com/cases/federal/us/596/20-303/}{\textit{United States v. Vaello Madero}, 596 U.S. 159, 180 (2022)}
\end{quote}

\section{Doctrinal Significance and Modern Legacy (Through 2026)}
\textit{Balzac} is widely regarded as the capstone of the \textit{Insular Cases}: it transformed Justice White's fractured \textit{Downes} concurrence into unanimous precedent and established that Congress's Territories Clause power allows it to hold populated territories indefinitely in an unincorporated status---even after conferring U.S. citizenship.

\begin{itemize}
  \item \textbf{Piecemeal Rights Extension vs. Plenary Benefits Discretion:} In the decades after \textit{Balzac}, the Court declined to extend the \textit{Insular Cases} to U.S. citizens tried by overseas courts-martial (\href{https://supreme.justia.com/cases/federal/us/354/1/}{\textit{Reid v. Covert}, 354 U.S. 1 (1957)}) and held that equal protection/due process (\href{https://supreme.justia.com/cases/federal/us/426/572/}{\textit{Examining Board v. Flores de Otero}, 426 U.S. 572 (1976)}) and the Fourth Amendment (\href{https://supreme.justia.com/cases/federal/us/442/465/}{\textit{Torres v. Puerto Rico}, 442 U.S. 465 (1979)}) apply in Puerto Rico. Yet on federal benefits, the Court relied on Territories Clause deference in \href{https://supreme.justia.com/cases/federal/us/446/651/}{\textit{Harris v. Rosario}, 446 U.S. 651 (1980)}, to uphold lower AFDC funding for Puerto Rico under rational-basis review.
  \item \textbf{\textit{Aurelius} (2020) and \textit{Vaello Madero} (2022):} In \href{https://supreme.justia.com/cases/federal/us/590/18-1334/}{\textit{Financial Oversight and Management Board for Puerto Rico v. Aurelius Investment, LLC}, 590 U.S. 448 (2020)}, the Court upheld PROMESA's oversight board as exercising ``primarily local'' powers without reaching whether to overrule the \textit{Insular Cases}. Two years later, in \href{https://www.law.cornell.edu/supremecourt/text/20-303}{\textit{United States v. Vaello Madero}, 596 U.S. 159 (2022)}, an 8--1 Court (per Justice Kavanaugh) held that excluding Puerto Rico residents from Supplemental Security Income (SSI) satisfied rational-basis review under the Fifth Amendment because residents are generally exempt from most federal income taxes. Crucially, however, \href{https://supreme.justia.com/cases/federal/us/596/20-303/}{Justice Gorsuch concurred} solely to urge overruling the \textit{Insular Cases}---explicitly indicting \textit{Balzac} and its predecessors as textually baseless relics rooted in racial stereotypes---while Justice Sotomayor dissented on equal-protection grounds and echoed calls to repudiate the doctrine.
  \item \textbf{Status Through 2026:} Despite bipartisan scholarly condemnation, a February 2024 \href{https://www.usccr.gov/files/2024-02/english_pr-ac_memo-1.pdf}{U.S. Commission on Civil Rights advisory memorandum} criticizing the unincorporated-territory doctrine, and Justice Gorsuch's invitation, the Supreme Court has not overruled \textit{Balzac}. In October 2022, the Court denied certiorari in \href{https://www.scotusblog.com/2022/10/court-declines-to-take-up-petition-seeking-to-overturn-insular-cases/}{\textit{Fitisemanu v. United States}} (leaving intact the Tenth Circuit's ruling that American Samoans are non-citizen U.S. nationals rather than birthright citizens under the Fourteenth Amendment). Through 2026, \textit{Balzac}'s incorporated/unincorporated divide remains formally on the books, even as its constitutional legitimacy is deeply contested.
\end{itemize}

\section{Pedagogical Critiques and Exam Hypotheticals}
\textbf{Pedagogical Critique:} \textit{Balzac} exposes a stark originalist and structural tension. Textually, the Territories Clause grants Congress power to make ``needful Rules and Regulations,'' not to suspend the Bill of Rights; as Justice Gorsuch noted in \textit{Vaello Madero}, the Constitution nowhere mentions ``incorporated'' versus ``unincorporated'' territories. Moreover, Taft's claim that jury trials are non-fundamental directly collides with \textit{Duncan v. Louisiana}, 391 U.S. 145 (1968), which later held the Sixth Amendment jury right ``fundamental to the American scheme of justice'' and incorporated it against the states---leaving \textit{Balzac}'s core premise doctrinally stranded.

\textbf{Exam Hypothetical (Standard Application):} Congress enacts an organic statute for a newly acquired, unincorporated Pacific territory that grants U.S. citizenship to all residents and includes a statutory Bill of Rights omitting grand-jury indictment. A resident citizen is prosecuted for a felony without a grand jury and challenges the conviction under the Fifth Amendment. \textit{Analysis:} Under \textit{Balzac}, unless Congress expressly declared incorporation (or manifested an implication ``so strong as to exclude any other view''), citizenship alone does not incorporate the territory, and non-fundamental procedural guarantees like grand-jury indictment do not apply \textit{ex proprio vigore}.

\textbf{``Opposite Side'' Contrast (Incorporated Territory or Fundamental Right):} Suppose instead that Congress enacts a statute expressly declaring a territory ``incorporated into the United States as a future State,'' or suppose the local government of an unincorporated territory authorizes warrantless, suspicionless searches of residents' homes or punishes political speech without due process. \textit{Contrast:} In the first scenario, express incorporation triggers the full Constitution---including Article~III and Sixth Amendment jury guarantees---just as in \textit{Rasmussen v. United States}. In the second scenario, even inside an unincorporated territory under \textit{Balzac}, ``fundamental personal rights'' such as due process and Fourth Amendment protections against unreasonable searches (\href{https://supreme.justia.com/cases/federal/us/442/465/}{\textit{Torres v. Puerto Rico}, 442 U.S. 465 (1979)}) apply directly to constrain territorial and federal action.

\end{document}
